Die vorliegende Simulation zielt darauf ab, die Herausforderungen, mit denen Rollstuhlnutzerinnen und Rollstuhlnutzer im urbanen Raum konfrontiert sind, so realitätsnah und nachvollziehbar wie möglich darzustellen. Neben der technischen Realisierung der Bewegungsplattform ist auch die Gestaltung der virtuellen Umgebung von entscheidender Bedeutung. Virtuelle Realitäten entfalten ihre Wirkung insbesondere dann, wenn Nutzende das Gefühl entwickeln, sich tatsächlich innerhalb der dargestellten Umgebung zu befinden. Dieses Empfinden wird in der Forschung häufig als Immersion bezeichnet (QUELLE)

Um eine solche immersive Erfahrung zu ermöglichen, müssen verschiedene gestalterische und konzeptionelle Aspekte berücksichtigt werden. Zu den relevanten Faktoren zählen demnach eine klare Nutzerführung, eine verständliche Informationsvermittlung, eine realitätsnahe audiovisuelle Gestaltung sowie eine strukturierte Präsentation der Inhalte innerhalb der Simulation. Im Rahmen des Projekts wurde daher mehrere Gestaltungselemente entwickelt, die darauf abzielen, die Orientierung der Nutzenden zu erleichtern und gleichzeitig ein möglichst realistisches Erlebnis der simulierten Umgebung zu ermöglichen. (QUELLE)

\section{Nutzerführung und Einführung in die Simulation}
\setauthor{Linnea Schönbauer}
Zu Beginn der Simulation erfolgt die Einführung der Nutzenden in die Anwendung über einen Startbildschirm. Der Startscreen ist ein elementarer Bestandteil der Simulation, dessen Funktionalität sich aus der Interaktion mit weiteren Elementen der Simulation ergibt. Einerseits dient er der kurzen inhaltlichen Einordnung des Projekts, indem er den Zweck der Simulation erläutert und den Nutzenden eine grundlegende Orientierung über den Ablauf vermittelt. Andererseits markiert der den Ausgangspunkt der eigentlichen Simulation. (QUELLE)

Die Integration eines solchen Einstiegsbildschirms ist insbesondere in interaktiven Anwendungen von Relevanz, da Nutzende ohne vorausgegangene Erklärung häufig Schwierigkeiten haben können, den Kontext der Simulation oder ihre eigene Roll innerhalb der virtuellen Umgebung zu verstehen. Die kurze textliche Einführung zu Beginn der Simulation dient dazu, den Nutzenden bereits vor dem Beginn der Anwendung einen Überblick über die zu erwartenden Erfahrungen zu verschaffen. (QUELLE)